\section{A self-contained bootstrap for a single moving factor}
\label{sec:diagonal-strip}

\paragraph{The diagonal test.}
Let $w_*$ denote the lower limit of $(1-\Re\rho)\log|\Im\rho|$
over the zeros of a fixed automorphic standard function as their height
tends to infinity. Suppose a prime-positive amplifier contains only this
function at moving nonscalar frequencies, with unit total target
weight, static pole coefficient $P$, and disk cost $G$.
Assume a preliminary zero-free estimate gives $w_*>0$ and that the
tilted disk comparison of the preceding note holds for this amplifier.
Then for every nonnegative normalized compact autocorrelation $h$,
\[
       w_* G\ge H(0)-P H(-1),\qquad
       H(z)=\int_0^\infty e^{-zt}h(t)\,dt.
\]
Indeed, if $w_*=\infty$ there is nothing to prove. Otherwise choose
zeros attaining the finite lower limit. For every $0<a<w_*$ there
is a strictly positive margin between $a$ and the scaled displacements
of all sufficiently high zeros. At height $\ell T+O(1)$ the ratio
of its logarithm to $\log|T|$ tends to one. Hence the same prior
holds throughout every required disk. Set $W=w_*$ in the tilted
comparison. It gives
$w_*G\ge H((w_*-a)/w_*)-P H(-a/w_*)$. The compact support makes
$H$ entire and allows $a\uparrow w_*$. Thus one uses the true lower
limit as a boundary; no sequence of revised constants is required.

\paragraph{An essentially self-dual representation.}
Let $K$ have degree $r$ and $\pi$ be a fixed essentially self-dual
unitary cuspidal representation of $\mathrm{GL}_n(\mathbb A_K)$.
Use the real trace $X$ after extracting a square root of the
essential twist locally, and the corresponding phase $y=\cos\phi$.
For $\alpha=1/\sqrt8$, the square
\[
          \alpha X^2+Xy+y^2/(4\alpha)
             =(2\alpha X+y)^2/(4\alpha)
\]
has static Rankin--Selberg coefficient $\alpha$, a moving standard
coefficient one, and static scalar coefficient $1/(8\alpha)$.
The scalar Fourier term at twice the height costs $o(\log|T|)$ by the
fixed-character Hecke estimate. Thus $P=1/\sqrt2$ and $G=nr/4$ by
convexity. Local factors at the omitted finite set of places contribute
$O(1)$, using their strict unitary exponent bounds and boundedness
of the compact weight. The fixed-rank theorem supplies $w_*>0$.

For $0<\theta<\pi/2$ put
\[
 P_0(\theta)=\frac{1-\theta\cot\theta}{\sin^2\theta}.
\]
Heath-Brown, \emph{Zero-free regions for Dirichlet $L$-functions},
Proc.~Lond.~Math.~Soc.~64 (1992), Lemmas~7.1 and~7.5 with scale one,
gives a nonnegative compact autocorrelation for which
$H(0)-P_0(\theta)H(-1)=P_0(\theta)\cos^2\theta$. This normalization
divides his function by its value at zero. Choosing
$\theta=1.3221231$ gives $P_0(\theta)>1/\sqrt2$ and
$P_0(\theta)\cos^2\theta>0.0428324150$.
Since $H(-1)\ge0$, the diagonal test therefore implies, under the
tilted disk comparison,
\[
        \min(\beta,1-\beta)>
          \frac1{5.83671\,nr\log|T|}
        \qquad(|T|\text{ sufficiently large}).
\]
The strict margin in the numerical comparison allows strict exclusion.
Functional-equation duality gives the other edge.

\paragraph{Elliptic power priors.}
For a fixed non-CM elliptic curve over $\mathbb Q$, use
$X=U_m(\cos\vartheta)$, the symmetric-power character, at a good
prime power. An amplifier requiring only this moving nonscalar can
use any additive majorants $A(X)+B(\cos\phi)$ of $|X\cos\phi|$.
Take $p=3/2$, $q=3$. Young's inequality bounds that product by
$t|X|^p/p+t^{-q/p}|\cos\phi|^q/q$ for $t>0$. Approximate each
continuous summand uniformly from above by an ordinary polynomial
on its compact range. On expanding in the $U_j$ characters and cosine
harmonics, all extra nonscalar characters are static. Their coefficients
may have either sign: each nontrivial fixed symmetric power is
holomorphic and nonzero at one. The poles come from the two constant
polynomial coefficients. These coefficients are the integrals against
Haar measure on $\mathrm{SU}(2)$ and on the circle, respectively,
by character and cosine orthogonality. This is an algebraic identification
of the constant term; no equidistribution theorem is necessary.
At the finitely many bad primes the fixed approximating polynomials
give bounded logarithmic derivatives in the weighted Euler sum.
The approximants are fixed first and the ordinate then tends to
infinity. Taking the infimum over $t$ shows that the pole can approach
arbitrarily closely from above
\[
 P_m=\left(\frac2\pi\int_0^\pi
    \left|\frac{\sin((m+1)v)}{\sin v}\right|^{3/2}
        \sin^2 v\,dv\right)^{2/3}
       \left(\frac4{3\pi}\right)^{1/3} .
\]
Here $4/(3\pi)$ is the circular mean of $|\cos\phi|^3$.
The coefficient of the target $U_m\cos\phi$ is exactly one. The
preliminary strip can be obtained from the preceding fixed-rank result.
Charge the target the Nelson exponent
$G=(n/4)(1-2/(3n^5-2n^4-n^2))$ with $n=m+1$.
One can let the arbitrarily small slack in the growth exponent tend
to zero after the height tends to infinity.

The following rows use strict upper budgets $P_m<P$ and angles with
$P<P_0(\theta)$. The last column rounds the certified lower limit
$P_0(\theta)\cos^2\theta/G$ downwards.
\[
\begin{array}{c|c|c|c}
 m&P&\theta&\text{usable scaled prior}\\\hline
 2&0.6881 &1.299692&0.0660323\\
 3&0.68308&1.293575&0.0512050\\
 5&0.6790 &1.28855 &0.0351170
\end{array}
\]
The interval calculation in \texttt{diagonal\_prior.py} encloses the
integral by $65536$ subintervals on $[0,\pi/2]$, using symmetry and
$|U_m|\le m+1$ on the first subinterval. On the others the sine
quotient is bounded by interval evaluation. For the original quartic
square one may thus take $a=0.051$ throughout. With $k=2.23970$,
$S=0.77992$, and $W=1/5.874730$, the positive scalar surplus exceeds
$2.061\cdot10^{-8}$. For the degree-twenty choice, $a=0.035$,
$k=2.32689$, $S=0.78016$, and $W=1/5.950712$, the surplus exceeds
$2.055\cdot10^{-8}$. The preceding comparison for $H$ gives both
elliptic consequences.

\paragraph{Verification scope.}
The tilted disk comparison, the analytic assembly with automorphic
functions and the uniform approximation argument for the static and
scalar majorants are human proofs. The enclosed constants and integrals
are independently checked by outward rounding. The diagonal limiting
implication and the pointwise Young inequality have direct Lean proofs.
